
\begin{register}{H}{INT\_ENABLE\_REG}{0x0104}\label{regdesc:INTENABLEREG}
\regfield{INT\_ENABLE}{31}{1}{{0x00000000}}%
\regfield{(reserved)}{1}{0}{{0}}%
\reglabel{Reset}\regnewline%
\begin{regdesc}
\begin{reglist}
\item [INT\_ENABLE\lbrack\regindex{n}\rbrack] Setting \regindex{n}th bit enables assertion of \regindex{n}th interrupt to the CPU. (R/W)
\begin{itemize}
    \setlength\itemsep{-1em}
    \item 0: Disabled;
    \item 1: Enabled;
\end{itemize}
\end{reglist}
\end{regdesc}
\end{register}

\begin{register}{H}{INT\_TYPE\_REG}{0x0108}\label{regdesc:INTTYPEREG}
\regfield{INT\_TYPE}{31}{1}{{0x00000000}}%
\regfield{(reserved)}{1}{0}{{0}}%
\reglabel{Reset}\regnewline%
\begin{regdesc}
\begin{reglist}
\item [INT\_TYPE\lbrack\regindex{n}\rbrack] Setting \regindex{n}th bit enables capturing the rising edge of \regindex{n}th interrupt. (R/W)
\begin{itemize}
    \setlength\itemsep{-1em}
    \item 0: Level type (signal level detection);
    \item 1: Pulse type (rising edge detection);
\end{itemize}
\end{reglist}
\end{regdesc}
\end{register}

\begin{register}{H}{INT\_CLEAR\_REG}{0x010C}\label{regdesc:INTCLEARREG}
\regfield{INT\_CLEAR}{31}{1}{{0x00000000}}%
\regfield{(reserved)}{1}{0}{{0}}%
\reglabel{Reset}\regnewline%
\begin{regdesc}
\begin{reglist}
\item [INT\_CLEAR\lbrack\regindex{n}\rbrack] Set \regindex{n}th bit to clear pending status of the \regindex{n}th interrupt. (R/W) \newline This is only useful for ``pulse" type interrupts, since ``level" type interrupts must be cleared at source. \newline Note that the set bit must be manually toggled back to 0 afterwards.
\begin{itemize}
    \setlength\itemsep{-1em}
    \item 0: Don't care;
    \item 1: Clear pending status;
\end{itemize}
\end{reglist}
\end{regdesc}
\end{register}

\begin{register}{H}{INT\_EIP\_REG}{0x0110}\label{regdesc:INTEIPREG}
\regfield{INT\_EIP}{31}{1}{{0x00000000}}%
\regfield{(reserved)}{1}{0}{{0}}%
\reglabel{Reset}\regnewline%
\begin{regdesc}
\begin{reglist}
\item [INT\_EIP\lbrack\regindex{n}\rbrack] Read \regindex{n}th bit to get the pending status of \regindex{n}th interrupt to CPU. (RO) \newline Only enabled and above threshold interrupts are reflected here.
\begin{itemize}
    \setlength\itemsep{-1em}
    \item 0: Not pending
    \item 1: Pending
\end{itemize}
\end{reglist}
\end{regdesc}
\end{register}

\begin{register}{H}{INT\_PRIORITY\_\regindex{n}\_REG (\regindex{n}: 1-31)}{0x0114+4*\regindex{n}}\label{regdesc:INTPRIORITYREG}
\regfield{(reserved)}{28}{4}{{0x00000000}}%
\regfield{INT\_PRIORITY\_\regindex{n}}{4}{0}{{0x0}}%
\reglabel{Reset}\regnewline%
\begin{regdesc}
\begin{reglist}
\label{fielddesc:INTPRIORITY}\item [INT\_PRIORITY\_\regindex{n}] Writing a 4-bit value to \regindex{n}th register configures priority of \regindex{n}th interrupt. (R/W) \newline \textit{Note : Interrupts with 0 priority are masked regardless of threshold value.}
\end{reglist}
\end{regdesc}
\end{register}

\begin{register}{H}{INT\_THRESH\_REG}{0x0194}\label{regdesc:INTTHRESHREG}
\regfield{(reserved)}{28}{4}{{0x00000000}}%
\regfield{INT\_THRESH}{4}{0}{{0x0}}%
\reglabel{Reset}\regnewline%
\begin{regdesc}
\begin{reglist}
\label{fielddesc:INTTHRESH}\item [INT\_THRESH] Writing a 4-bit value configures the global priority threshold for all interrupts. (R/W) \newline All interrupts with priority lower than the threshold are masked. \newline \textit{Note : Interrupts with 0 priority are masked regardless of threshold value.}
\end{reglist}
\end{regdesc}
\end{register}
